\scanpage{191}
\mspage{743}

structure -- liée à une classe de sous-groupes discrets commensurables entre eux -- \unsure{des|de} \deleted{\unsure{méta-}}\inserted{$L$-$T$-}groupes. (où $L$ est initiale de ``lacets'', et $T$ de ``Teichmüller'') -- ce sont des sortes de \unsure{mixtures} de groupes de Teichmüller et de $L$-groupes. Pour illustrer ce point, partons d'une extension
\diagram{(6)\qquad
\begin{tikzcd}[ampersand replacement=\&, column sep=large]
1 \arrow[r] \& \Pi \arrow[r] \& \mathcal{G} \arrow[r] \& \mathcal{T} \arrow[r] \& 1
\end{tikzcd}
}
où $\Pi$ est muni d'une $L$-structure invariante par $\mathcal{G}$, alors $\mathcal{G}$ est un $L$-$T$-groupe admettant (si $\Pi$ est $\overline{\mathrm m}$ un \deleted{\unsure{$L$-}}groupe à lacets) $\Pi$ comme sous-groupe de définition. Partant du groupe à lacets $\Pi$, on peut prendre p.ex.
\[
(7)\qquad
\left\{
\begin{array}{l}
\mathcal{G}=\mathrm{Aut}_{L}(\Pi)\\
\mathcal{T}=\mathrm{Autext}_{L}(\Pi)
\end{array}
\right.
\qquad\text{d'où}
\,,
\]

\scanpage{192}

et ainsi le groupe \unsure{ambiant} $\mathcal{G}$, extension du groupe de Teichmüller $\mathcal{T}$ par le groupe à lacets $\Pi$, apparaît comme un $LT$-groupe.

\deleted{\unsure{2}}

\scanpage{193}
\mspage{744}

2$^\circ$) \uline{Cas profini}

L'extension de la notion de $L$-groupe ou de $LT$-groupe \inserted{$\Pi$}, en termes de \uline{germes} d'hom extérieurs
\diagram{(8)\qquad
\begin{tikzcd}[ampersand replacement=\&]
\unsure{K_i^{\mathcal{G}}|\kappa_i^{\mathcal{G}}}\colon T \arrow[r] \& \Pi
\end{tikzcd}\,,
}
est essentiellement évidente. Voici l'exemple le plus important pour nous. \deleted{\unsure{Soit}}

Soit d'abord $k$ un corps \deleted{alg. clos de car. $0$}, et soit $M$ une \deleted{\unsure{0-conn.}} \deleted{\unsure{0-connexe}} multiplicité algébrique sur $k$, qui soit \inserted{0-connexe,} de t.f. sur $k$, \deleted{lisse} séparée$/k$ et lisse \unsure{en|de} dim relative $1$. On doit pouvoir \unsure{prouver} alors que $M$ peut se décrire par une courbe alg. lisse et quasi-projective $U$ sur $k$, un groupe fini \deleted{$\unsure{K}$} $G$ qui opère sur $U$, \deleted{et} comme
\[
(9)\qquad M=(U,G)=
\begin{array}{l}
\inserted{\text{(petit)}}\text{ Topos annelé des}\\
\text{faisceaux étales sur}\\
U\text{ où }G\text{ opère}\ldots
\end{array}
\]

\scanpage{194}

\deleted{\unsure{ou}} -- Plus intrinsèquement, il doit être vrai qu'il existe un rev. étale fini de $M$ (qu'on peut supposer galoisien, de groupe $G$) \deleted{\unsure{tel|et}} qui \unsure{soit|est} \unsure{une|un} \deleted{courbe algébrique}\inserted{schéma} -- laquelle \unsure{sera|ser.} néc. \unsure{quasi-projectif|quasi-projective}, lisse, de dim relative $1$ sur $k$. Supposons en tous cas qu'il en soit ainsi, choisissons un pt base géom. $a$ de $M$, et considérons
\[
(10)\qquad \Pi=\Pi_1(M,a)\simeq\Pi_1(\unsure{U|M},G;a)\,.
\]

Nous supposerons désormais $k$ de car. zéro. Supposons d'abord $k$ alg clos. Utilisant la suite exacte
\diagram{(11)\qquad
\begin{tikzcd}[ampersand replacement=\&, column sep=large]
1 \arrow[r] \& \underbrace{\Pi_1(U,a)}_{\Pi_0} \arrow[r] \& \underbrace{\Pi_1(\unsure{U|M},G;a)}_{\Pi} \arrow[r] \& G \arrow[r] \& 1
\end{tikzcd}
}
et la \deleted{\unsure{$L$}} structure à lacets naturelle sur $\Pi_0$, on trouve une $L$-structure sur $\Pi=\Pi_1(\unsure{M},a)$, qui manifestement ne

\scanpage{195}
\mspage{745}

dépend pas du choix du rev. étale \inserted{fini} galoisien $U$ de $M$, et du relèvement choisi du pt géom. $a$ de $M$ en un pt. géom. de $U$. Nous munirons tjrs \inserted{tacitement} $\Pi=\Pi_1(M,a)$ de cette $L$-structure.

Si à nouveau $k$ est un corps \unsure{quelcon-que|quelconque} (le cas le plus intéressant étant celui où $k$ est ext. de type fini de $\mathbb{Q}$), on aura une suite exacte
\diagram{(12)\qquad
\begin{tikzcd}[ampersand replacement=\&, column sep=large]
1 \arrow[r] \& \Pi_1(M_{\overline{k}},a) \arrow[r] \& \Pi_1(M,a) \arrow[r] \& \Gamma_{\overline{k}/k} \arrow[r] \& 1
\end{tikzcd}
}
et $\Pi_1(M_{\overline{k}},\unsure{a|o})$ apparaît comme muni d'une $L$-structure invariante par $\Pi_1(M,\unsure{a|o})=\Pi$, de sorte que $\Pi$ apparaît comme un $TL$-groupe.

Je suis particulièrement intéressé au cas ``universel'' de type $(g,\nu)$ -- mais la base \inserted{naturelle} n'est pas alors le

\scanpage{196}
\mspage{\deleted{\unsure{745}}}

spectre d'un corps $k$, mais la \deleted{\unsure{schéma}} \inserted{multiplicité} modulaire $M_{g,\nu}$. [À vrai dire, pour la structure de $L$-groupe discret $\Pi_0$, jouant le rôle de groupe de référence, on \deleted{peut} \inserted{doit pouvoir} introduire une \deleted{\unsure{schéma}} \inserted{multiplicité} modulaire correspondante sur $\mathbb{Q}$, et $\overline{\mathrm m}$ sur $\mathbb{Z}$. Le fait qui risque de faire des difficultés \unsure{conceptuelles} pour une telle construction, c'est qu'une multiplicité \inserted{``modulaire''} (sur $\mathbb{Q}$, ou sur $\mathbb{Z}$\ldots) décrit en principe une \uline{catégorie} fibrée --- telle celle des schémas en courbes relatifs de type $(g,\nu)$ --- alors que les $\unsure{1|L}$-multiplicités relatives correspondantes à un ``type'' $\Pi_0$-donné forment une \uline{2-catégorie} \inserted{(en groupoïdes)} fibrée sur $(\mathrm{Sch})$, et pas seulement une \uline{catégorie} fibrée\ldots]

\scanpage{197}
\mspage{746}

On aura donc une suite exacte
\diagram{\begin{tikzcd}[ampersand replacement=\&, column sep=small]
(13)\quad 1 \arrow[r] \&
\widehat{\Pi}_{g,\nu} \arrow[r] \&
\underbrace{\Pi_1({\unsure{U|M}_{g,\nu}}_{\mathbb{Q}})}_{\mathcal{M}_{g,\nu}} \arrow[r] \&
\underbrace{\Pi_1({M_{g,\nu}}_{\mathbb{Q}})}_{\mathcal{N}_{g,\nu}} \arrow[r] \& 1,
\end{tikzcd}}
où ${U_{g,\nu}}_{\mathbb{Q}}$ est la courbe \unsure{de|du} type $g,\nu$ universelle'' (en car. $0$). Ici c'est $\Pi_1(U_{g,\nu}\unsure{\deleted{\mathbb{Q}}|\xi}_{\mathbb{Q}})$ qui apparaît comme un $TL$-groupe profini.

Je voudrais dégager une propriété commune à ce cas ``universel'' (13), \deleted{\unsure{et}} \deleted{(qui \unsure{se} \unsure{ramène} d'ailleurs \unsure{essentielle-|essentiellement}} \inserted{et au cas d'une multiplicité} à un corps de base $k$ extension \uline{de type fini} de $\mathbb{Q}$ (le premier cas se ramenant d'ailleurs essentiellement au second, en prenant un ``point générique'' de $M_{g,\nu}$\ldots). Pour ceci, considérons un $LT$-groupe $\unsure{\mathcal{M}|\mathcal{A}}$, et un \deleted{de} homomorphisme
\diagram{\begin{tikzcd}[ampersand replacement=\&]
\unsure{k_i|\kappa_i}:T \arrow[r] \& \deleted{\unsure{\mathcal{G}|\Pi}}\,\mathcal{M}
\end{tikzcd}}
appartenant à un des $\unsure{K_i^{\mathcal{G}}|\kappa_i^{\mathcal{G}}}$. On peut

\scanpage{198}

p. ex. choisir un sous-groupe de définition distingué $\Pi$ dans $\mathcal{M}$, qui sera donc un groupe à lacets, \deleted{\unsure{et on aura|et on a}} une suite exacte
\diagram{\begin{tikzcd}[ampersand replacement=\&, column sep=small]
(14)\quad 1 \arrow[r] \& \Pi \arrow[r] \& \mathcal{M} \arrow[r] \& \mathcal{N} \arrow[r] \& 1,
\end{tikzcd}}
et un \unsure{hom|lacet}
\diagram{(15)\qquad \begin{tikzcd}[ampersand replacement=\&]
\unsure{k|\kappa}:T \arrow[r,hook] \& \Pi
\end{tikzcd},\qquad
\unsure{k|\kappa}(T)\overset{\mathrm{dfn}}{=}L
}
relatif à la structure à lacets de $\Pi$. Considérons le normalisateur de $\unsure{L|T}$ dans $\mathcal{M}$
\[
(16)\qquad \operatorname{Norm}_{\mathcal{M}}(L)\overset{\mathrm{dfn}}{=}N(L)
\]
et considérons l'action de $N(L)/L\subset\mathcal{N}$ \deleted{\unsure{dans}} sur $L$. (\uline{NB} On voit aisément que l'image de $N(L)/L$ dans $\mathcal{N}$ est un sous-groupe ouvert de $\mathcal{N}$\ldots) Cette action se fait par un ``\unsure{caractère} cyclotomique'', induit par un \unsure{caractère}
\diagram{(17)\qquad \begin{tikzcd}[ampersand replacement=\&]
\mathcal{N} \arrow[r,"\chi_{\mathcal{N}}"] \& \widehat{\mathbb{Z}}^{*}
\end{tikzcd}.}
L'hypothèse essentielle est que l'image

\scanpage{199}
\mspage{747}

de $\deleted{\chi}\;\mathcal{N}$ par $\chi$ (ou aussi de $N(L)/L$ par $\chi$) soit un sous-groupe ouvert de $\widehat{\mathbb{Z}}^{*}$. Cette hypothèse a un caractère intrinsèque au $LT$-groupe $\mathcal{M}$ lui-$\overline{\mathrm m}$ (indépendamment du choix de $\Pi$) --- car le caractère cyclotomique
\diagram{(18)\qquad \begin{tikzcd}[ampersand replacement=\&]
\chi_{\mathcal{M}}:\mathcal{M} \arrow[r] \& \widehat{\mathbb{Z}}^{*}
\end{tikzcd}}
a un caractère intrinsèque à $\mathcal{M}$, et l'hypothèse signifie que l'image de ce dernier dans $\widehat{\mathbb{Z}}^{*}$ est ouverte.

Revenant à $L$ et $N(L)$, on voit que l'hypothèse \unsure{faite|faîte} implique ceci:
\[
\begin{array}{ll}
(19)&\text{Tout sous-groupe ouvert de }N(L)/L\\
&\text{opère non trivialement sur }L.
\end{array}
\]

On sait, par des arguments \unsure{standarts|standards} que ceci implique que pour tout \unsure{bon|\deleted{bon}} corps $K$ \deleted{de car. $\unsure{0}$} (notamment un corps $\ell$-adique $\mathbb{Q}_{\ell}$),

\scanpage{200}

et tout hom. de $\mathcal{M}$ dans un groupe $G(K)$; où $G$ est un schéma en groupes de t.f. affine sur $K$; \deleted{\unsure{le|u}}
\diagram{(20)\qquad \begin{tikzcd}[ampersand replacement=\&]
u:\mathcal{M} \arrow[r] \& G(K)
\end{tikzcd},}
la restriction de $u$ à $L$ est ``quasi-unipotente'', i.e.
\[
\begin{array}{ll}
(21)&\exists\ \text{sous-groupe ouvert }L'\text{ de }L,\text{ tel que}\\
&u(L')\subset G(K)\text{ soit formé d'éléments}\\
&\text{unipotents},
\end{array}
\]
--- ce qui revient au $\overline{\mathrm m}$, si $\ell_i$ est un générateur de $L=L_i$, qu'on ait
\[
\begin{array}{ll}
(22)&u(\ell_i)\text{ est un élément quasi-}\\
&\text{unipotent de }G(K),\text{ i.e. }\exists\ n\in\mathbb{N}^{*},\\
&\text{tel que }\ell_i^{n}\text{ soit unipotent}.
\end{array}
\]

\scanpage{201}
\mspage{748}

\textcircled{\(3^\circ\)}\quad \uline{Enveloppe schématique d'un} \(T\)\unsure{-}\(L\)-\uline{groupe profini} : \inserted{\unsure{\(a)\)}} \uline{Les hom. admissibles.}

Soient \(\mathcal{M}\) un \(T\)-\(L\)-groupe profini, dont le caractère structural (18)
\diagram{\begin{tikzcd}[ampersand replacement=\&]
\mathcal{M} \arrow[r, "\chi_{\mathcal{M}}"] \& \widehat{\mathbb{Z}}^{*}
\end{tikzcd}}
ait une image ouverte, \(\Pi\) un sous-groupe de définition \unsure{distingué} (c'est donc un \deleted{\unsure{\(L\)-}}groupe \inserted{à lacets} profini\unsure{)|]}, et \(\Pi_0\subset\Pi\) un sous-groupe \uline{discret} de \(\Pi\), qui en soit une ``discrétification'' --- c'est donc un \unsure{groupe} à lacets discret, \unsure{etc.|et} \(\Pi\) s'identifie à son complété profini\ldots

Soit \(G_{\mathbb{Q}}\) un schéma en groupes \inserted{affine de t.f.} sur \(\mathbb{Q}\). Soit \(\mathcal{A}_{\mathbb{Q}}\) l'anneau des adèles de \(\mathbb{Q}\)
\[
(23)\qquad \mathcal{A}_{\mathbb{Q}}=\widehat{\mathbb{Z}}\otimes_{\mathbb{Z}}\mathbb{Q}\subset\prod_{\ell\in\mathbb{P}}\mathbb{Q}_{\ell}
\]
et considérons \unsure{\deleted{le}|le} groupe des points adéliques
\[
(24)\qquad G(\mathcal{A}_{\mathbb{Q}})\subset\prod_{\ell\in\mathbb{P}}G(\mathbb{Q}_{\ell}).
\]
Si \(G_{\mathbb{Q}}\) provient de un \unsure{schém.|schéma} en groupes

\scanpage{202}

affine de t.f. \(G\) sur \(\mathbb{Z}\), alors on a
\[
\begin{gathered}
(25)\qquad G(\mathcal{A}_{\mathbb{Q}})\simeq
\unsure{(|\{}(x_{\ell})_{\ell\in\mathbb{P}}\in\prod_{\ell}G(\mathbb{Q}_{\ell})
\mid x_{\ell}\in G(\mathbb{Z}_{\ell})\\
\hfill\text{pour presque }\unsure{\forall|\mathrm{tt}}\ell\}
\end{gathered}
\]

\deleted{On dit qu'\unsure{un|on}}

On se donne de plus un hom
\diagram{\begin{tikzcd}[ampersand replacement=\&]
(26)\quad G_{\mathbb{Q}} \arrow[r, "\chi_G"] \& \unsure{\mathbb{G}_{m\mathbb{Q}}|G_{m\mathbb{Q}}}.
\end{tikzcd}}
On s'intéresse à \deleted{\unsure{un|ces}}\inserted{des} homomorphismes de groupes
\diagram{\begin{tikzcd}[ampersand replacement=\&]
(27)\quad \mathcal{M} \arrow[r] \& G(\mathcal{A}_{\mathbb{Q}})
\end{tikzcd}}
satisfaisant à des propriétés que nous allons énumérer.

\(1^\circ)\) \uline{Continuité}. Cela équivaut à dire que \(a)\) les hom \(\mathcal{M}\longrightarrow G(\unsure{\mathbb{Q}_{\ell}})\) sont \unsure{continus|continues}, et \(b)\) \(\exists S\) partie finie de \(\mathbb{P}\), telle que pour \(\ell\notin S\), on ait \deleted{\illegible{1}} une factorisation (continue)
\diagram{\begin{tikzcd}[ampersand replacement=\&]
(28)\quad \mathcal{M} \arrow[r] \& G(\widehat{\mathbb{Z}}_{\ell}) \& (\ell\in\mathbb{P}\setminus S)
\end{tikzcd}}
--- condition qui en effet ne dépend pas du ``modèle sur \(\mathbb{Z}\)'' \(G\) de \(G_{\mathbb{Q}}\) choisi.

\uline{NB} (28) est un hom. de groupes profinis.

\scanpage{203}
\mspage{749}

\uline{ReNB} Sauf erreur, la propriété \(1^\circ)\) signifie que l'on peut trouver un schéma en groupes \(G\) affine de t.f. sur \(\mathbb{Z}\), qui \unsure{définisse|définit} \(G_{\mathbb{Q}}\), et tel que (27) se factorise par \(G(\widehat{\mathbb{Z}})\).
\diagram{\begin{tikzcd}[ampersand replacement=\&]
(29)\quad \deleted{\unsure{G}}\quad \mathcal{M} \arrow[r] \& G(\widehat{\mathbb{Z}}).
\end{tikzcd}}

\(2^\circ)\) \uline{Compatibilité avec} \(\chi\) :

\(\unsure{(30)}\)
\diagram{\begin{tikzcd}[ampersand replacement=\&, row sep=large, column sep=large]
\& \mathcal{M} \arrow[r] \arrow[dl, "\chi_{\mathcal{M}}"'] \& G(\mathcal{A}_{\mathbb{Q}}) \arrow[dl, "\chi_G"] \\
\widehat{\mathbb{Z}} \arrow[r, hook] \& \unsure{\mathbb{G}_m|G_m}(\mathcal{A}_{\mathbb{Q}})\simeq\mathcal{A}_{\mathbb{Q}}^{*}\subset\displaystyle\prod_{\ell\in\mathbb{P}}\mathbb{Q}_{\ell}^{*} \&
\end{tikzcd}}
est commutatif.

\(3^\circ)\). \uline{Compatibilité avec la discrétification} \(\Pi_0\).

\((31)\)
\diagram{\begin{tikzcd}[ampersand replacement=\&, row sep=large, column sep=large]
\Pi'_0 \arrow[r, hook] \arrow[d, "\text{factorisation}"] \& \Pi \arrow[r, hook] \& \mathcal{M} \arrow[d] \\
G(\mathbb{Q}) \arrow[rr, hook] \&\& G(\mathcal{A}_{\mathbb{Q}})
\end{tikzcd}},
pour un sous-groupe d'indice fini convenable \(\Pi'_0\) de \(\Pi_0\).

\uline{NB} Quitte à changer de notations \unsure{ops|qps} \(\Pi'_0=\Pi_0\).

\(4^\circ)\) \uline{Caractère minimal} (\unsure{irrédondance|irrédundance}) :

Si \(G'_{\mathbb{Q}}\) est un sous-schéma en groupes de

\scanpage{204}

\(G_{\mathbb{Q}}\), tel que \(\deleted{\unsure{G}}\ \mathcal{M}\longrightarrow G(\mathcal{A}_{\mathbb{Q}})\) se factorise par \(G'(\mathcal{A}_{\mathbb{Q}})\), alors \(G=G'\).

\uline{NB} Si cette hypothèse n'est pas vérifiée, alors parmi tous les \(G'_{\mathbb{Q}}\subset G_{\mathbb{Q}}\) qui la mettent en défaut, il y en a un \inserted{(unique)} plus petit, et en remplaçant \(\chi_G\) par \(\chi_{G'}=\chi_G\mid G'_{\mathbb{Q}}\ldots\), les hyp. \(1^\circ):4^\circ)\) sont satisfaites pour \(G'_{\mathbb{Q}},\chi'\), et \(\mathcal{M}\longrightarrow G'(\mathcal{A}_{\mathbb{Q}})\).

\uline{Remarque}. Soit \(\mathcal{G}_0\subset\mathcal{M}\) le sous-groupe de \(\mathcal{M}\) formé des \(g\in\mathcal{M}\) qui sont tels que \(\unsure{g(\Pi_0)|g\ell(\Pi_0)}\) et \(\Pi_0\) soient commensurables. \unsure{Càd.|C.à.d.} \inserted{les \(g\)} qui normalisent le ``germe de discrétification'' \(\Pi_0\) de \(\Pi\).

\margin{\(\mathcal{G}_0/\unsure{\Pi_0}\) joue le rôle d'un groupe de Teichmüller \uline{discret}}

Il est possible qu'il faille renforcer la condition \(3^\circ\), en exigeant que l'on ait
\diagram{\begin{tikzcd}[ampersand replacement=\&]
\mathcal{G}_0 \arrow[r] \& G(\mathbb{Q}).
\end{tikzcd}}
L'avenir le décidera !

\scanpage{205}
\mspage{750}

Réflexion \unsure{faite|fait}, je préfère poser les choses de façon légèrement \unsure{différente|différent}, en introduisant, dans le \(LT\)-groupe \inserted{profini} \(\mathcal{M}\), un sous-groupe \deleted{\unsure{définissant la \(LT\)-structure|définissant la \(LT\)-struct.}} distingué \(\Pi\) avec une \(L\)-structure (pas \unsure{néc.\textsuperscript{t}|néc.} une structure à lacets --- p.ex. on peut avoir \(\Pi=\mathrm{SL}\unsure{(}2,\mathbb{Z})\)) définissant la \(LT\)-structure de \(\mathcal{M}\), une discrétification \(\Pi_0\) de \(\Pi\) (dans un sens évident, pour les \(L\)-structures) un sous-groupe discret \(\mathcal{G}_0\) \inserted{\(\supset\Pi_0\)} normalisant \(\Pi_0\) (NB. en cas particuliers \unsure{importants} on \(\mathcal{G}_0=\Pi_0\)), tel que \inserted{son adh.} \(\unsure{\widehat{\mathcal{G}}_0|\mathcal{G}_0^{*}}\) \inserted{\(\overset{\mathrm{dfn}}{=}\mathcal{G}\)} dans \(\mathcal{M}\) soit distingué dans \(\mathcal{M}\).

On a donc un diagramme de groupes

\((32)\)
\diagram{\begin{tikzcd}[ampersand replacement=\&, 
row sep=large, column sep=large,
execute at end picture={
  \path (s205p.north) -- coordinate[midway] (s205pgmid) (s205g.north);
  \path (s205g.north) -- coordinate[midway] (s205gmmid) (s205m.north);
  \path (s205p.north) -- coordinate[midway] (s205pmmid) (s205m.north);
  \path (s205pzero.south) -- coordinate[midway] (s205zeromid) (s205gzero.south);
  \draw ([yshift=5mm]s205p.north) .. controls ([yshift=8mm]s205p.north) and ([xshift=-2mm,yshift=5mm]s205pgmid) .. ([yshift=8mm]s205pgmid) .. controls ([xshift=2mm,yshift=5mm]s205pgmid) and ([yshift=8mm]s205g.north) .. ([yshift=5mm]s205g.north);
  \node[above] at ([yshift=8mm]s205pgmid) {$\mathcal{T}$};
  \draw ([yshift=5mm]s205g.north) .. controls ([yshift=8mm]s205g.north) and ([xshift=-2mm,yshift=5mm]s205gmmid) .. ([yshift=8mm]s205gmmid) .. controls ([xshift=2mm,yshift=5mm]s205gmmid) and ([yshift=8mm]s205m.north) .. ([yshift=5mm]s205m.north);
  \node[above] at ([yshift=8mm]s205gmmid) {$\Gamma$};
  \draw ([yshift=17mm]s205p.north) .. controls ([yshift=20mm]s205p.north) and ([xshift=-2mm,yshift=17mm]s205pmmid) .. ([yshift=20mm]s205pmmid) .. controls ([xshift=2mm,yshift=17mm]s205pmmid) and ([yshift=20mm]s205m.north) .. ([yshift=17mm]s205m.north);
  \node[above] at ([yshift=20mm]s205pmmid) {$\mathcal{N}$};
  \draw ([yshift=-3mm]s205pzero.south) .. controls ([yshift=-6mm]s205pzero.south) and ([xshift=-2mm,yshift=-3mm]s205zeromid) .. ([yshift=-6mm]s205zeromid) .. controls ([xshift=2mm,yshift=-3mm]s205zeromid) and ([yshift=-6mm]s205gzero.south) .. ([yshift=-3mm]s205gzero.south);
  \node[below] at ([yshift=-6mm]s205zeromid) {$\mathcal{T}_0$};
}]
\unsure{1} \arrow[r, hook] \& |[alias=s205p]| \underset{=\widehat{\Pi}_0}{\Pi} \arrow[r, hook] \& |[alias=s205g]| \underset{=\widehat{\mathcal{G}}_0}{\unsure{\widehat{\mathcal{G}}|\mathcal{G}^{*}}_{\unsure{\deleted{0}|0}}} \arrow[r, hook] \& |[alias=s205m]| \mathcal{M} \\
1 \arrow[r, hook] \& |[alias=s205pzero]| \Pi_0 \arrow[u, hook] \arrow[r, hook] \& |[alias=s205gzero]| \mathcal{G}_0 \arrow[u, hook] \&
\end{tikzcd}},

où la première ligne est formée de groupes profinis, la deuxième de groupes discrets. On notera que \(\mathcal{G}\) s'identifie à un quotient du complété profini de \(\mathcal{G}_0\), mais

\scanpage{206}

je n'exige pas pour le moment qu'il soit isomorphe à ce complété lui-$\overline{\mathrm m}$. \unsure{Donc} \unsure{idem|item} pour $\mathcal{T}$ et $\mathcal{T}_0$. Le groupe $\mathcal{T}_0\inserted{=\mathcal{G}_0/\Pi_0}$ joue le rôle d'un groupe de Teichmüller \inserted{``absolu''} discret, le groupe $\mathcal{T}$ celui d'un complété de gr. de Teichmüller, le groupe $\Gamma=\mathcal{M}/\mathcal{G}$ joue le rôle d'un groupe de Galois profini. \deleted{La \unsure{situation}} Je suis intéressé surtout par deux situations extrêmes

a) $\mathcal{G}_0=\Pi_0$ d'où $\mathcal{G}=\Pi$ --- situation provenant d'une \unsure{$1$-multiplicité} lisse relative sur un corps \deleted{$k$} ext. de t.f. de $\mathbb{Q}$, ou plus généralement \deleted{d'un schéma}\inserted{sur un} schéma \inserted{(ou une multiplicité)} \unsure{int.} de type fini sur $\mathbb{Q}$, \deleted{\unsure{$\Gamma$|$k$} \unsure{en fait}}\inserted{$\mathcal{N}\unsure{\xrightarrow{\sim}|\simeq}\Gamma$, est} \deleted{\unsure{est}} le groupe fondamental \deleted{de cette} schéma ou de la dite multiplicité,

b) $\mathcal{G}_0/\Pi_0\simeq$ groupe de Teichmüller de $\Pi_0$ \inserted{ou d'indice fini dans ce gr. de T.} --- dans ce cas, on est intéressé surtout à une situation ``modulaire'', où

\scanpage{207}
\mspage{751}

$\mathcal{N}$ s'interprète comme le groupe fondamental de \unsure{la|leur} multiplicité modulaire universelle, $\mathcal{T}$ comme sa partie ``géométrique'', $\mathcal{T}_0$ comme sa partie ``transcendante'', et $\Gamma$ comme $\Gamma_{\unsure{\overline{\mathbb{Q}}_0|\overline{\mathbb{Q}}}/\mathbb{Q}}$ ....

On fait sur (32) l'hypothèse que le caractère cyclotomique sur $\mathcal{M}$ (opérant sur $\Pi$ par autom. de $L$-structures) soit trivial sur $\mathcal{G}$, a fortiori sur $\Pi$ (donc
\[
(33)\qquad \Pi=\Pi^{+},\quad \mathcal{G}=\mathcal{G}^{+}\qquad ),
\]
i.e. qu'il se factorise en
\diagram{(34)\qquad
\begin{tikzcd}[ampersand replacement=\&, column sep=large]
\Gamma \arrow[r, "\chi_{\Gamma}"] \& \widehat{\mathbb{Z}}^{*}
\end{tikzcd},
}
\uline{dont l'image est un ss-groupe ouvert.}

Si $(G_{\mathbb{Q}},\chi_G)$ est comme plus haut (p. 748), on s'intéresse aux hom (27)
\diagram{
\begin{tikzcd}[ampersand replacement=\&, column sep=large]
\mathcal{M} \arrow[r] \& G(\mathcal{A}_{\mathbb{Q}})
\end{tikzcd}
}
satisfaisant les conditions $1^{\circ})$ $2^{\circ})$ $\deleted{\unsure{2}}\inserted{3}^{\circ})$ $4^{\circ})$ comme dessus, mais où $3^{\circ})$ est \unsure{formulée} sous la forme un peu plus forte

($3^{\circ}$ forte) (Compatibilité avec discrétification de Teichmüller $\mathcal{G}_0$) : On a une \unsure{factorisa-}

\scanpage{208}

tion $\mathcal{G}_0\longrightarrow G(\mathbb{Q})$, rendant \unsure{donc} commutatif le diagramme
\diagram{(35)\qquad
\begin{tikzcd}[ampersand replacement=\&, column sep=large, row sep=large]
\mathcal{G}_0 \arrow[r, hook] \arrow[d] \& \mathcal{M} \arrow[d] \\
G(\mathbb{Q}) \arrow[r, hook] \& G(\mathcal{A}_{\mathbb{Q}})
\end{tikzcd}.
}

Les structures de la forme (32) méritent sans doute un nom, je vais \unsure{les} appeler --- ne serait-ce que provisoirement, \uline{\unsure{structures} de Galois-Teichmüller.}
\margin{Quitte par la suite à introduire des conditions plus draconiennes sur \unsure{cette structure}}

Si $G$ est un groupe algébrique affine sur $\mathbb{Q}$, un hom $\mathcal{M}\unsure{\longrightarrow|{-}}G(\mathcal{A}_{\mathbb{Q}})$ satisfaisant les conditions $1^{\circ})$ $2^{\circ})$ $3^{\circ})$ \inserted{(forte)} $4^{\circ})$ est appelé un ``hom de Lie'' de la structure de Galois-Teichmüller.

Je vais désigner par la suite par $\underline{\mathcal{M}}_{\mathbb{Q}}$ \inserted{(ou \unsure{simplement} $\underline{\mathcal{M}}$ ?)} le schéma en groupes \inserted{$G_{\mathbb{Q}}$} sur $\mathbb{Q}$. \deleted{\unsure{défini}}

Je désigne par \deleted{$\unsure{\mathcal{G}|\mathcal{M}}$} $\underline{\mathcal{G}}_{\mathbb{Q}}$ le sous-schéma en groupes de $\underline{\mathcal{M}}_{\mathbb{Q}}$ engendré par l'

\scanpage{209}
\mspage{752}

image de $\mathcal{G}_0\longrightarrow\deleted{\unsure{\mathcal{M}|G}}G(\mathbb{Q})$, par $\underline{\Pi}_{\mathbb{Q}}$ celui engendré par l'image de $\Pi\longrightarrow G(\mathbb{Q})$

On a ainsi \deleted{\unsure{un}} des sous-schémas en groupes
\[
(36)\qquad \underline{\Pi}_{\mathbb{Q}}\subset
\deleted{\underline{\unsure{\mathcal{M}}}}\inserted{\underline{\mathcal{G}}}_{\mathbb{Q}}
\subset\underline{\mathcal{M}}_{\mathbb{Q}}
\]
\uline{distingués} dans $\underline{\mathcal{M}}_{\mathbb{Q}}$, donnant lieu à des schémas en groupes quotients
\[
(37)\qquad
\left\{
\begin{aligned}
&\underline{\mathcal{T}}_{\mathbb{Q}}\simeq
\underline{\mathcal{G}}_{\mathbb{Q}}/\underline{\Pi}_{\mathbb{Q}},\quad
\underline{\Gamma}_{\mathbb{Q}}\simeq
\underline{\mathcal{M}}_{\mathbb{Q}}/\underline{\mathcal{G}}_{\mathbb{Q}},\\
&\qquad\underline{\mathcal{N}}_{\mathbb{Q}}\simeq
\underline{\mathcal{M}}_{\mathbb{Q}}/\underline{\Pi}_{\mathbb{Q}}
\end{aligned}
\right.
\]
--- donnant lieu à quatre suites exactes de schémas en groupes, dont la dernière est
\diagram{(38)\qquad
\begin{tikzcd}[ampersand replacement=\&, column sep=large]
1 \arrow[r] \& \underline{\mathcal{T}}_{\mathbb{Q}} \arrow[r] \&
\underline{\mathcal{N}}_{\mathbb{Q}} \arrow[r] \& \underline{\Gamma}_{\mathbb{Q}} \arrow[r] \& 1
\end{tikzcd}.
}

\deleted{Je \unsure{vais|vois} \unsure{réfléchir|rappeler}}\inserted{Notons \unsure{que}} le caractère $\chi_G:G_{\mathbb{Q}}\longrightarrow \unsure{\mathbb{G}_m|G_m}$ est épi, et se factorise par $\underline{\Gamma}_{\mathbb{Q}}$ en
\diagram{(39)\qquad
\begin{tikzcd}[ampersand replacement=\&, column sep=large]
\underline{\Gamma}_{\mathbb{Q}} \arrow[r, "\chi_{\underline{\Gamma}}"] \& \unsure{\mathbb{G}_{m\mathbb{Q}}|G_{m\mathbb{Q}}}
\end{tikzcd}.
}
